% Required for PDF/A metadata.
% Since LaTeX 2025-11-01, \DocumentMetadata always loads the
% latex-lab tagging modules. This template only needs PDF management
% for the PDF/A metadata, not the experimental tagging layer, so use
% the dedicated pdfmanagement package when it is available.
\IfFileExists{pdfmanagement.sty}{%
  \RequirePackage{pdfmanagement}%
  \SetKeys[document/metadata]{pdfstandard = a-2b, pdfversion = 1.7, lang = sl}%
}{%
  % Compatibility fallback for older TeX installations.
  \providecommand\DebugTemplatesOn{}%
  \providecommand\DebugTemplatesOff{}%
  \DocumentMetadata{pdfstandard = a-2b, pdfversion = 1.7, lang = sl}%
}

% if row numbering is required by your supervisor for review, use the lineno option
\documentclass[9pt]{style/friteza}
% use the following file to define/include additional elements
\graphicspath{{fig/}}

% ------------ additional commands/macros/includes defined by the user
\etal{in sod.}
\renewcommand{\Authands}{ in }
\renewcommand{\Authand}{ in }

\newcommand{\set}[1]{\ensuremath{\mathbf{#1}}}
\renewcommand{\vec}[1]{\ensuremath{\mathbf{#1}}}
\newcommand{\uvec}[1]{\ensuremath{\hat{\vec{#1}}}}
\newcommand{\const}[1]{{\ensuremath{\kappa_\mathrm{#1}}}}

\usepackage{subcaption}
\usepackage{csquotes}
\usepackage{svg}
\usepackage{silence}
\WarningFilter{latex}{Command \showhyphens has changed}

\newmintinline[ltx]{latex}{breaklines}
% --------------------------------------------------

\title[slovene]{Primerjava implementacij oddaljenega nadgrajevanja procesorjev ESP32-S3 in NRF5340}
\title[english]{OTA implementation comparison between ESP32-S3 and NRF5340}

\author{Arne Salobir}
% Please give the surname of the lead author for the running footer
\leadauthor{Arne S.}

\supervisor{prof. dr.}{Patricio Bulić}
% If you do not have a cosupervisor, comment out the following row.
%\cosupervisor{prof. dr.}{Jure Demšar}
%\secondcosupervisor{prof. dr.}{Primer Drugega Somentorja}

% Choose your study programmes (and track, if applicable): BUN-RI, BVS-RI, BUN-RM, BMA-RI-PV, BMA-RI, BMA-MM, BMA-UI
\program{BUN-RI}
% Study-year
\vol{2025/26}
% Choose main language: slovene or english
\mainlanguage{slovene}

\keywords[slovene]{ESP32, nRF5340, OTA, BLE, Wi-Fi}

\keywords[english]{}

% Provide a link (or links) to the repository where your data and code is published. If no code was produced, comment out this section.
\repository{\href{https://github.com/Drsoldier/espOtaDiploma}}

\sloabstract{}

\sloabstract{Mikrokrmilniki so danes prisotni v številnih napravah, od pametnih luči in avtomobilskih sistemov do industrijskih aplikacij. Vendar programsko opremo potrebujemo včasih tudi spremeniti. V tem primeru lahko vsak mikrokrmilnik posamezno nadgradimo ali pa uporabimo brezžične tehnologije za nadgrajevanje programske opreme s postopkom OTA (Over-The-Air). Namen diplomskega dela je primerjava implementacij brezžičnih tehnologij Wi-Fi in Bluetooth Low Energy (BLE) za nadgrajevanje programske opreme in katera tehnologija je učinkovitejša za implementacijo in uporabo. Za primerjavo implementacij smo izbrali platformo ESP32-S3 za Wi-Fi in nRF5340-DK za BLE. Pri tem smo za programiranje ESP32-S3 uporabljali orodje Espressif-IDF, katerega je razvil Espressif za ESP32 mikrokrmilnike. Za programiranje nRF5340-DK smo uporabili nRF Connect SDK, katerega je razvil Nordic Semiconductor za njihove mikrokrmilnike. Na podlagi implementacij smo primerjali katero brezžično tehnologijo je lažje implementirati ter tudi napisali skripte, katere se uporabljajo za testiranje le-teh.}

\engabstract{TODO}

\begin{document}
\maketitle

% If your first paragraph (i.e. with the \dropcap) contains a list environment (quote, quotation, theorem, definition, enumerate, itemize...), the line after the list may have some extra indentation. If this is the case, add \parshape=0 to the end of the list environment.
\dropcap{D}iplomsko delo primerja implementaciji brezžičnih tehnologij Wi-Fi in BLE za nadgrajevanje vgrajene programske kode preko brezžičnih omrežji (Firmware Over-The-Air) za mikrokrmilnika ESP32-S3 in NRF5340-DK. Nadgrajevanje vgrajene programske kode ima veliko različnih pristopov, ter posledično veliko različnih rešitev. Manjše število mikrokrmilnikov in večja potreba po nizki porabi energije je bolj primerna za uporabo BLE pri nadgrajevanju kode. Če večja poraba energije ni pomembna za naše odločitve, ali da imamo veliko več mikrokrmilnikov, je bolje izbrati Wi-Fi.  

Za mikrokrmilnike z dostopom do interneta in drugih brezžičnih tehnologij, poznano kot IoT (Internet-of-Things) poznamo že od leta 2001 dalje \cite{embeddedEverywhere}. Vendar z tako ravzpitimi in hitrim razvojom IoT naprav se znajdejo velike težave v kodi +\cite{secureFOTAForIoT}, ki so včasih zelo kritične in potrebujejo drugačno kodo. Pred razvojem brezžičnih tehnologij, bi podjetja potrebovala vsak mikrokrmilnik predelati nazaj, medtem ko OTA omogoča brezžično modificiranje že obstoječe kode. 

Za pravilno programiranje mikrokrmilnikov potrebujemo komplet za razvijanje programske opreme (SDK), saj je vsak procesor narejen drugače. V tem delu se bomo omejli na mikrokrmilnike ESP32-S3 proizvajalca Espressif in NRF5340 proizvajalca Nordic Semiconductor ter njune SDK:  Espressif-IDF\cite{espressifMain} in nRF Connect SDK\cite{nordicSemiSDK}. 

\section{Metode}

\section{Rezultati}


\section{Sklep}

% The acknowledgemenst content (not mandatory).
\acknow{Raziskavo je omogočila Javna agencija za raziskovalno dejavnost Republike Slovenije (ARRS) v okviru programa P2-0359, Vseprisotno računalništvo.}

% Bibliography
\bibliography{{bib/bibliography}}

\end{document}
